% CHAPTER 2: BACKGROUND STUDY AND LITERATURE REVIEW
\proposalchapter{Background Study and Literature Review}

\section{Background Study}

\subsection{Study of Existing System}

\textbf{Manual system.} In the traditional system, a passenger books a long-distance bus ticket by visiting the operator's counter, calling the office, or going through a local agent. The staff check a paper register or a seat chart, tell the passenger which seats are free, take the payment in cash, and write the booking by hand. Each counter or agent keeps its own record, so seat information is not shared in real time. This leads to double booking, long queues, errors in the records, and difficulty in producing sales reports.

\textbf{Existing online platforms.} Several online platforms already offer bus ticket booking in Nepal. BusSewa describes itself as a real-time booking platform that lets passengers search routes, compare fares, choose seats, and receive e-tickets, with payment through wallets, banks, and cards~\cite{bussewa}. GoMyGo allows several passengers to be booked in one transaction, sends confirmations by email and SMS, and follows each operator's own cancellation policy~\cite{gomygo}. Okay Journey offers an interactive seat map and a choice of bus class~\cite{okayjourney}. Trip Turbo lets users filter buses by price and departure time, pick window or aisle seats from a seat map, and receive an e-ticket by email and SMS~\cite{tripturbo}.

\textbf{Comparison.} Table~\ref{tab:comparison} compares the manual system, the existing online platforms, and the proposed system.

\begin{table}[htbp]
\centering
\caption{Comparison of the manual system, existing platforms, and the proposed system}
\label{tab:comparison}
\small
\singlespacing
\renewcommand{\arraystretch}{1.25}
\newcolumntype{L}[1]{>{\raggedright\arraybackslash}p{#1}}
\begin{tabular}{|L{2.7cm}|L{3.1cm}|L{3.5cm}|L{3.5cm}|}
\hline
\textbf{Feature} & \textbf{Manual system} & \textbf{Existing online platforms} & \textbf{Proposed system}\\
\hline
Seat availability & Checked at the counter or by phone & Real-time & Real-time\\
\hline
Seat selection & Done by staff, limited & Interactive seat map & Interactive seat map\\
\hline
Booking record & Paper registers & Stored in a database & Stored in a database\\
\hline
Payment & Cash at the counter & Wallets, banks, cards & Simulated (prototype)\\
\hline
Notifications & None & Email and SMS & Not included\\
\hline
Operator reports & Hard to prepare & Offered through the platform & Bookings, revenue, occupancy\\
\hline
Scope & Single counter & Commercial, many operators & Academic prototype, sample data\\
\hline
\end{tabular}
\end{table}

\textbf{Summary.} Existing platforms have solved the main problems of the manual system for the operators they work with. This project does not aim to replace them. Its purpose is to build the core booking logic as a complete system with its own modules: seat locking to prevent double booking, fare rules, sorting and searching of results, role-based management for passengers, operators, and administrators, and reports. The modules are written for this project and not taken from plugins or third-party services, so they can be studied, tested, and extended.

\subsection{Key Concepts}
[To be written]

\section{Literature Review}
[To be written]